\begin{definition}[Independent virtual coordinates of a region]%
    \label{def:peps_injective_vertex_coordinates}
    \lean{TNLean.PEPS.RegionVertexVirtualConfig,
        TNLean.PEPS.regionVertexTensorMap,
        TNLean.PEPS.regionVertexLeftInverse,
        TNLean.PEPS.regionIncidentVertexConfig,
        TNLean.PEPS.regionVirtualBondWeight,
        TNLean.PEPS.regionVirtualBondMap}
    \leanok
    Give each vertex of $R$ an independent configuration of its incident
    virtual labels. Thus each internal edge has two labels, one at each
    endpoint. Let $F_R=\bigotimes_{v\in R}A_v$ be the product site map.
    For injective sites let $L_R$ be the product of their left inverses.
    Let $B_R$ denote the virtual region map that pairs the two endpoint
    labels of every internal edge and leaves the crossing labels open.
    Its coefficients are sums of basis vectors indexed by the actual
    incident-edge assignments, with the boundary labels fixed.
\end{definition}

\begin{theorem}[Exact virtual reconstruction of a region]%
    \label{thm:peps_injective_vertex_reconstruction}
    \lean{TNLean.PEPS.regionVertexLeftInverse_comp_regionVertexTensorMap,
        TNLean.PEPS.regionVertexTensorMap_injective,
        TNLean.PEPS.regionVertexTensorMap_single,
        TNLean.PEPS.regionVertexTensorMap_regionVirtualBondWeight,
        TNLean.PEPS.openRegionMap_eq_regionVertexTensorMap_comp_regionVirtualBondMap,
        TNLean.PEPS.regionVertexLeftInverse_comp_openRegionMap,
        TNLean.PEPS.regionGroundSpace_eq_map_regionVirtualBondMap_range,
        TNLean.PEPS.comap_regionGroundSpace_regionVertexTensorMap}
    \leanok
    \uses{def:peps_injective_vertex_coordinates,def:peps_region_ground_space}
    For arbitrary site tensors the actual open-region map satisfies
    $T_R=F_RB_R$, and $\mathcal G_R=F_R(\operatorname{ran}B_R)$.
    If all site maps are injective, then
    \[
        L_RF_R=I,\qquad L_RT_R=B_R,\qquad
        F_R^{-1}(\mathcal G_R)=\operatorname{ran}B_R.
    \]
    No nonzero-bond hypothesis is needed for these identities.
    These are the finite-graph virtual-coordinate identities underlying
    the parent argument in \cite[Section IV.C.1]{Cirac2021Matrix}.
\end{theorem}

\begin{proof}\leanok
    The product site map takes a virtual basis vector to the product of
    its site tensor coefficients. Summing over compatible internal edge
    labels with the crossing labels fixed gives the actual open-region
    contraction, so $T_R=F_RB_R$. Tensoring the site left inverses gives
    $L_RF_R=I$, hence injectivity of $F_R$ and $L_RT_R=B_R$.
    The range and inverse-image formulas follow.
\end{proof}
